% ============================================================
% SECTION 9: AGE OF ANALYSIS (~1700 - 1800)
% Slides 66-70
% ============================================================

% ---- Section Divider ----
{
\setbeamercolor{background canvas}{bg=enlightenmentTeal}
\begin{frame}[plain]
  \centering
  \vfill
  \textcolor{white}{\Huge\bfseries The Age of Giants}\\[0.7em]
  \textcolor{white!85}{\Large When Mathematics Conquered the Physical World}\\[1.5em]
  \textcolor{white!65}{\large 1700--1800 CE}
  \vfill
  \visualplaceholder[5cm]{Diagram}
  \vfill
\end{frame}
}

\section{The Age of Analysis}

% ---- Slide 67: Euler ----
\begin{frame}{Leonhard Euler --- The Most Prolific Mathematician in History}
  \begin{columns}[T]
    \column{0.52\textwidth}
    \begin{personbox}[Leonhard Euler]{enlightenmentTeal}
      \textbf{Dates:} 1707 -- 1783 CE\\
      \textbf{Origin:} Basel, Switzerland;\\worked in St.\ Petersburg and Berlin
    \end{personbox}
    \vspace{0.3em}
    \begin{itemize}\setlength\itemsep{2pt}
      \item Published $\sim$866 papers and books --- more than any other individual mathematician
      \item Standardized notation: $e$, $i$, $\pi$, $f(x)$, $\Sigma$
      \item Founded \textbf{graph theory} (K\"{o}nigsberg bridges)
      \item Foundational contributions to number theory, topology, mechanics, optics, astronomy
      \item Continued producing mathematics \emph{after going blind}
    \end{itemize}

    \column{0.45\textwidth}
    % Euler's identity
    \visualplaceholder[5cm]{Five constants labeled}

    \vspace{0.3em}
    % Konigsberg bridges
    \visualplaceholder[5cm]{Four land masses}
  \end{columns}

  \note{
    \textbf{Talking point:} ``He was so productive that even after going blind, he continued
    to produce mathematics by dictating to scribes.''

    The K\"{o}nigsberg bridge problem (1736) is considered the origin of both graph theory
    and topology. Euler proved that no path crossing each bridge exactly once exists
    by analyzing the structure of the connections (what we now call vertex degrees).

    \verifytag{Publication count: sometimes cited as $\sim$800--900. The Euler Archive
    lists 866 works. Say ``approximately 866'' or ``over 800.''}
  }
\end{frame}

% ---- Slide 68: Lagrange ----
\begin{frame}{Joseph-Louis Lagrange \hfill {\small\textcolor{enlightenmentTeal!60}{[EXTENDED]}}}
  \begin{columns}[T]
    \column{0.52\textwidth}
    \begin{personbox}[Joseph-Louis Lagrange]{enlightenmentTeal}
      {\small (born Giuseppe Lodovico Lagrangia)}\\[3pt]
      \textbf{Dates:} 1736 -- 1813 CE\\
      \textbf{Origin:} Turin, Sardinia (modern Italy);\\
      worked in Berlin and Paris
    \end{personbox}
    \vspace{0.3em}
    \begin{itemize}
      \item \textit{M\'{e}canique Analytique} (1788) --- reformulated Newtonian mechanics using \textbf{pure analysis}
      \item Boasted: ``no diagrams in this book''
      \item \textbf{Lagrange multipliers} --- optimization with constraints
      \item \textbf{Lagrangian mechanics} --- foundation of modern theoretical physics
      \item Major contributions to number theory
    \end{itemize}

    \column{0.45\textwidth}
    % Euler-Lagrange equation
    \visualplaceholder[5cm]{Diagram}

    \vspace{0.5em}
    \visualplaceholder[4cm]{Portrait of Lagrange}
  \end{columns}

  \note{
    Lagrange represents the triumph of analysis: he believed that all of mechanics could
    be reduced to pure calculus, without geometric diagrams. His \textit{M\'{e}canique Analytique}
    was the first systematic treatise to do this.

    For the audience: Lagrangian mechanics is what physics students learn in university.
    It is more powerful than Newton's force-based approach for complex systems.
  }
\end{frame}

% ---- Slide 69: Laplace ----
\begin{frame}{Pierre-Simon Laplace \hfill {\small\textcolor{enlightenmentTeal!60}{[EXTENDED]}}}
  \begin{columns}[T]
    \column{0.52\textwidth}
    \begin{personbox}[Pierre-Simon Laplace]{enlightenmentTeal}
      {\small (marquis de Laplace)}\\[3pt]
      \textbf{Dates:} 1749 -- 1827 CE\\
      \textbf{Origin:} Beaumont-en-Auge, Normandy, France
    \end{personbox}
    \vspace{0.3em}
    \begin{itemize}
      \item \textit{M\'{e}canique C\'{e}leste} --- applied calculus to celestial mechanics, completing Newton's program
      \item \textbf{Laplace transform} --- essential tool in engineering and physics
      \item Major contributions to \textbf{probability theory} (Bayesian reasoning, central limit theorem precursors)
    \end{itemize}

    \vspace{0.3em}
    {\footnotesize When Napoleon asked where God was in his system:}\\
    {\small\itshape ``Sire, I had no need of that hypothesis.''}

    \column{0.45\textwidth}
    % Laplace equation
    \visualplaceholder[5cm]{Diagram}

    \vspace{0.5em}
    \visualplaceholder[4cm]{Portrait of Laplace}
  \end{columns}

  \note{
    Laplace essentially completed Newton's program of explaining the solar system
    mathematically. Where Newton had left loose ends (perturbations of planetary orbits),
    Laplace showed that the solar system was stable.

    The Napoleon anecdote is famous but may be apocryphal --- or at least simplified.
    It captures Laplace's approach: mathematics explains the universe without needing
    supernatural intervention.
  }
\end{frame}

% ---- Slide 70: Gauss ----
\begin{frame}{Carl Friedrich Gauss --- \textit{Princeps Mathematicorum}}
  \begin{columns}[T]
    \column{0.50\textwidth}
    \begin{personbox}[Carl Friedrich Gauss]{enlightenmentTeal}
      \textbf{Dates:} 1777 -- 1855 CE\\
      \textbf{Origin:} Braunschweig (Brunswick), Germany;\\
      worked in G\"{o}ttingen
    \end{personbox}
    \vspace{0.3em}
    \begin{itemize}\setlength\itemsep{2pt}
      \item \textit{Disquisitiones Arithmeticae} (1801) --- foundational work in number theory
      \item Proved the \textbf{fundamental theorem of algebra}
      \item Discovered \textbf{non-Euclidean geometry} (unpublished!)
      \item The \textbf{Gaussian distribution} (bell curve)
      \item First major result at age 19: constructibility of the regular \textbf{17-gon}
    \end{itemize}

    \column{0.47\textwidth}
    % Gaussian bell curve (pure TikZ)
    \visualplaceholder[5cm]{Bell curve (filled)}

    \vspace{0.3em}
    % 17-gon
    \visualplaceholder[5cm]{Regular 17-gon}
  \end{columns}

  \vspace{0.1em}
  \centering
  {\small\itshape ``Gauss bridges two centuries. The 19th century would bring a crisis --- and a transformation --- in what mathematics IS\ldots''}

  \note{
    \textbf{Talking point:} ``He was called the `Prince of Mathematicians' and kept notebooks
    full of results he never published. Other mathematicians kept rediscovering things
    Gauss had already found.''

    At age 19, Gauss proved that the regular 17-gon can be constructed with compass and
    straightedge --- a problem open since antiquity. He was so excited he switched from
    studying languages to mathematics.

    Gauss spans two centuries: born 1777, died 1855. He anchors the Age of Analysis but
    his later work influenced 19th-century mathematics deeply.
  }
\end{frame}
